\chapter{Experiments and results}
\label{ch:results}

The experiments were run in the order of Table~\ref{tab:results-timeline}, and each answered a
question raised by the one before. This chapter is organised by question instead. Most settings
were trained once, with seed 42. Training noise turned out to be large
(Section~\ref{sec:results-seeds}), so the main comparison, with and without the TI-RADS-motivated
features, was repeated with three seeds, and single-run differences below about 0.03--0.05 AUC are
not interpreted.

\section{Timeline of experiments}
\label{sec:results-timeline}

\begin{table}[ht]
  \centering
  \small
  \caption{Experiments. Records are folders in \texttt{runs/}; the consolidated report is
           \texttt{reports/consolidated\_report\_2026-10-12.md}.}
  \label{tab:results-timeline}
  \begin{tabular}{lp{4.3cm}p{6.2cm}}
    \toprule
    date & experiment & outcome \\
    \midrule
    2 Oct  & Triage layer, run 1: mixed perturbations (official split, v1, 8 features, Qwen2.5-1.5B) & nothing certified; perturbations push the model off its training format \\
    2 Oct  & Triage layer, run 2: format-preserving perturbations & nothing certified; votes no better than softmax; calibration and official test not exchangeable \\
    7 Oct  & Error analysis and miss-rate variant (no new training) & errors are stable across queries; miss rate certifiable at 5\% for few nodules \\
    8 Oct  & Backbone sweep on the official split (Qwen2.5-1.5B retrained, Qwen2.5-7B) & 7B not better \\
    8 Oct  & Pooled split; image quality control; XGBoost & exchangeability restored; feature problems found \\
    8--9 Oct & v2 features; LLM runs A--F on the pooled split & v2 helps (+0.05 AUC); format and size results \\
    9 Oct  & Llama-2-7B with the settings of Ra et al.\ (run G); rich format (run H) & neither helps \\
    9--10 Oct & TI-RADS-motivated features (runs I, J, K) & largest gain of all changes; seventh feature and rich format add nothing \\
    11 Oct & Qwen2.5-7B at a lower learning rate, with TI-RADS features (run L) & clearly worse than 1.5B \\
    12 Oct & Seed repeats with and without TI-RADS features (runs M--P); fusion with XGBoost; conformal checks & TI-RADS gain confirmed; the language model adds nothing to XGBoost; Mondrian coverage holds over re-splits \\
    \bottomrule
  \end{tabular}
\end{table}

\section{Learning from radiomics prompts}
\label{sec:results-llm}

Yes, to a moderate degree. On the official split, with eight v1 features, Qwen2.5-1.5B reached
test AUC 0.672 (validation 0.704, calibration 0.707), against 0.648--0.663 for a logistic
regression, a random forest and a support vector machine. A retrained copy of the same model gave
0.667 on test, and its scores agreed closely with the first model (Spearman correlation 0.97).
Validation and calibration AUCs above the test AUC are consistent with the shift described in
Section~\ref{sec:data-official}.

Table~\ref{tab:results-llm} lists the runs on the pooled split. With the radiomics features alone
the best settings reach a test AUC of 0.75--0.76; with the TI-RADS-motivated features, 0.785 on
average over three seeds.

\begin{table}[ht]
  \centering
  \small
  \setlength{\tabcolsep}{4pt}
  \caption{Language-model runs on the pooled split (16 radiomics features, 5 epochs; learning rate
           $2 \times 10^{-4}$, LoRA on Q/V and seed 42 unless stated). Sensitivity at threshold
           0.5 on test. (f): near-duplicate filter; +T6, +T7: six or seven TI-RADS-motivated
           features. Run G uses the settings of Ra et al.\ (learning rate $10^{-4}$, LoRA on all
           linear layers); run L uses learning rate $5 \times 10^{-5}$.}
  \label{tab:results-llm}
  \begin{tabular}{llllrrrrr}
    \toprule
    run & model & features & format & val & cal & test & sens. & minutes \\
    \midrule
    A & Qwen2.5-1.5B & v2            & decimal    & 0.753 & 0.757 & 0.750 & 0.66 & 26 \\
    B & Qwen2.5-1.5B & v1 (f)        & decimal    & 0.709 & 0.705 & 0.699 & 0.67 & 27 \\
    C & Qwen2.5-7B   & v2 (f)        & decimal    & 0.708 & 0.705 & 0.713 & 0.00 & 90 \\
    D & Qwen2.5-1.5B & v2 (f)        & decimal    & 0.750 & 0.752 & 0.754 & 0.71 & 26 \\
    E & Qwen2.5-1.5B & v2 (f)        & z-score    & 0.729 & 0.756 & 0.759 & 0.65 & 37 \\
    F & Qwen2.5-1.5B & v2 (f)        & percentile & 0.709 & 0.713 & 0.734 & 0.65 & 37 \\
    G & Llama-2-7B   & v2 (f)        & decimal    & 0.705 & 0.709 & 0.731 & 0.66 & 182 \\
    H & Qwen2.5-1.5B & v2 (f)        & rich       & 0.732 & 0.711 & 0.744 & 0.64 &  83 \\
    I & Qwen2.5-1.5B & v2 (f) +T6    & decimal    & 0.795 & 0.812 & 0.804 & 0.65 & 36 \\
    J & Qwen2.5-1.5B & v2 (f) +T7    & decimal    & 0.786 & 0.802 & 0.781 & 0.39 & 37 \\
    K & Qwen2.5-1.5B & v2 (f) +T6    & rich       & 0.803 & 0.785 & 0.801 & 0.64 & 132 \\
    L & Qwen2.5-7B   & v2 (f) +T6    & decimal    & 0.751 & 0.712 & 0.693 & 0.56 & 123 \\
    \midrule
    M & I, seed 1    & v2 (f) +T6    & decimal    & 0.771 & 0.777 & 0.768 & 0.66 & 37 \\
    N & I, seed 2    & v2 (f) +T6    & decimal    & 0.787 & 0.794 & 0.783 & 0.69 & 37 \\
    O & D, seed 1    & v2 (f)        & decimal    & 0.706 & 0.695 & 0.712 & 0.69 & 36 \\
    P & D, seed 2    & v2 (f)        & decimal    & 0.723 & 0.715 & 0.713 & 0.65 & 36 \\
    \bottomrule
  \end{tabular}
\end{table}

Training was unstable in several runs. In the early epochs validation sensitivity at 0.5 was
sometimes close to 0 (all nodules called benign) and recovered later. Selecting the epoch by
validation AUC hides this from the final AUC, but not from threshold-based measures such as
sensitivity (run J ends at 0.39).

\section{How large is the training noise?}
\label{sec:results-seeds}

Larger than the effects of most changes. The best setting (run I) was trained again with seeds 1
and 2 (runs M and N), and the reference without TI-RADS features (run D) likewise (runs O and P).
Only the initialisation of the adapters and the order of the training batches change; the split
and the selected features stay the same.

\begin{itemize}
  \item The three seeds of the best setting give test AUC 0.804, 0.768 and 0.783: mean 0.785,
        standard deviation 0.018. Seed 42 was the upper end. Runs M and I differ significantly on
        every split (test interval $[-0.058, -0.014]$).
  \item The three seeds without TI-RADS features give 0.754, 0.712 and 0.713; again seed 42 was
        the best.
  \item A repeat of run H with the same seed, where only the padding length differed, moved by
        0.046 on test.
\end{itemize}

A paired bootstrap interval holds the trained model fixed and resamples nodules, so it does not
contain this variation. Comparisons between single runs therefore understate the uncertainty, and
the result that matters most for the thesis, the effect of the TI-RADS-motivated features, is
reported over matched seeds (Section~\ref{sec:results-clinical}).

\section{Do better features help?}
\label{sec:results-features}

Yes, clearly. Runs B and D differ only in the feature set: the test AUC rises from 0.699 with v1
to 0.754 with v2. The paired bootstrap interval for run A against run B is $[+0.013, +0.089]$ on
test and $[+0.015, +0.090]$ on calibration, so the gain of about 0.05 is not a chance
fluctuation of the test set.

XGBoost shows the same picture (Table~\ref{tab:results-xgb}). With v1, adding features helps
slowly; with v2, 16 features already give 0.78--0.81.

\begin{table}[ht]
  \centering
  \caption{XGBoost test AUC on the pooled split, by number of mRMR features. ``Filtered'':
           near-duplicates ($|r| > 0.95$) removed first; the last column gives the number of
           features left.}
  \label{tab:results-xgb}
  \begin{tabular}{lrrrr}
    \toprule
    features & 8 & 16 & 32 & all \\
    \midrule
    v1           & 0.695 & 0.723 & 0.738 & 0.793 (67) \\
    v1, filtered & 0.691 & 0.717 & 0.743 & 0.797 (44) \\
    v2           & 0.723 & 0.780 & 0.803 & 0.800 (102) \\
    v2, filtered & 0.723 & \textbf{0.805} & 0.796 & 0.800 (54) \\
    \bottomrule
  \end{tabular}
\end{table}

The near-duplicate filter lifts XGBoost with 16 v2 features from 0.780 to 0.805, because the
16 slots are no longer spent on copies. For the language model it makes no difference (runs A and
D, within 0.005).

On the official split, more features had seemed useless: classical models gained 0.05--0.10 AUC
on validation by going from 8 to 67 features but not on the official test set. On the pooled split
the gain carries over to the test set, so the earlier result was a consequence of the shift
between the official sources, not of the features.

\section{Does the number format matter?}
\label{sec:results-format}

Saying what the number is does not matter; how extreme the number is does; and adding what the
value means for the nodule does not help.

\begin{table}[ht]
  \centering
  \caption{AUC differences between formats, 95\% paired bootstrap intervals.}
  \label{tab:results-format}
  \begin{tabular}{llrr}
    \toprule
    comparison & what changes & calibration & test \\
    \midrule
    E $-$ D & z-score wording               & $[-0.016, +0.022]$ & $[-0.017, +0.029]$ \\
    F $-$ D & percentile values             & $[-0.075, -0.005]$ & $[-0.057, +0.022]$ \\
    H $-$ D & rich format                   & $[-0.074, -0.011]$ & $[-0.041, +0.023]$ \\
    K $-$ I & rich format, with TI-RADS     & $[-0.049, -0.005]$ & $[-0.031, +0.022]$ \\
    \bottomrule
  \end{tabular}
\end{table}

\begin{itemize}
  \item The z-score wording (run E) and the bare decimal (run D) perform the same.
  \item Percentiles (run F) lose about 0.04 AUC on validation and calibration; on test the
        interval includes zero. A percentile maps a z-score of 2 and a z-score of 5 to nearly the
        same value (98 and 100), so the model loses the information that a value is extreme.
        The loss suggests that the model uses the magnitude of the numbers, not only their order.
  \item The \emph{rich} format (run H) keeps the z-score and adds the percentile, a band
        (``lower than typical'') and what the value means (``meaning a less round, more irregular
        outline''). It does not help: test AUC 0.744 against 0.754, and 0.04 lower on calibration,
        where the interval excludes zero. A repeat with the same seed gave 0.698.
  \item With the TI-RADS-motivated features (run K against run I), the rich format is equal on
        validation and test and slightly worse on calibration. Its prompts are up to three times
        as long, and training took 132 minutes instead of 36. Training was stable, so the weak
        rich runs without these features were not caused by prompt length alone.
\end{itemize}

Plain-language meaning gives the model nothing it does not already get from the numbers. The
decimal format is kept.

\section{Does a larger backbone help?}
\label{sec:results-size}

No: neither a larger Qwen model, at two learning rates, nor the backbone and settings of Ra et al.\
beat Qwen2.5-1.5B.

\begin{itemize}
  \item On the official split, Qwen2.5-7B matched the 1.5B model: test AUC 0.669 with 8 features
        and 0.666 with 16, against 0.667 for 1.5B.
  \item On the pooled split, Qwen2.5-7B (run C) was worse than 1.5B (run D): 0.713 against 0.754,
        paired interval $[-0.069, -0.004]$ on test.
  \item With the backbone and settings of Ra et al.\ (run G: Llama-2-7B, learning rate
        $10^{-4}$, LoRA on all linear layers), the test AUC is 0.731, again below run D. The
        paired intervals against run D are $[-0.070, -0.016]$ on calibration and
        $[-0.052, +0.007]$ on test; against Qwen2.5-7B (run C) they include zero.
  \item With the TI-RADS-motivated features and a four times lower learning rate (run L),
        Qwen2.5-7B reaches 0.693, far below the 1.5B model with the same features (run I, 0.804;
        intervals exclude zero on every split).
\end{itemize}

The first 7B training did not settle. Across the five epochs, validation sensitivity at threshold
0.5 went 0.00, 0.01, 0.93, 0.04, 0.62; gradient norms reached 200; the best epoch by validation AUC
was epoch 2, and the selected model calls every test nodule benign at 0.5 although its ranking
still has AUC 0.713. The learning rate of $2 \times 10^{-4}$, chosen for the 1.5B model, was the
obvious suspect, so run L used $5 \times 10^{-5}$. This calmed the swings (sensitivity 0.85, 0.30,
0.31, 0.46, 0.58) but not the result: gradient norms still reached 100--200, validation AUC crept
up to 0.753 while the learning rate decayed to zero, and calibration and test fell well below
validation (0.712 and 0.693). The TI-RADS features, which lift the 1.5B model by about 0.06, gave
Qwen2.5-7B nothing (run L against run C: test interval $[-0.062, +0.023]$). Run G took 182
minutes on one L4 GPU, seven times as long as Qwen2.5-1.5B, for a lower AUC.

Within this setup (QLoRA on Q and V, about 2,000 training prompts, 5 epochs), a larger backbone
does not help, and Qwen2.5-1.5B stays the main model.

\section{Do TI-RADS-motivated features help?}
\label{sec:results-clinical}

Yes, for both models, and more than any other change tried. Table~\ref{tab:results-clinical-uni}
shows each feature on its own, on the training split.

\begin{table}[ht]
  \centering
  \small
  \caption{TI-RADS-motivated features on their own: AUC on the training split (above 0.5: higher
           in malignant nodules) and the direction TI-RADS leads one to expect.}
  \label{tab:results-clinical-uni}
  \begin{tabular}{llrl}
    \toprule
    category & feature & train AUC & as expected? \\
    \midrule
    shape          & taller-than-wide ratio & 0.658 & yes (higher in cancers) \\
    margin         & solidity               & 0.367 & yes (lower in cancers) \\
    composition    & anechoic fraction      & 0.409 & yes (benign nodules more cystic) \\
    echogenicity   & echogenicity ratio     & 0.552 & no (expected below 0.5) \\
    margin         & margin sharpness       & 0.563 & no (expected below 0.5) \\
    echogenic foci & punctate foci density  & 0.487 & no signal \\
    \bottomrule
  \end{tabular}
\end{table}

Three features behave as TI-RADS predicts, and the taller-than-wide ratio alone is as
discriminative as the best single radiomics features. The echogenicity ratio points the wrong way:
it correlates $-0.78$ with the anechoic fraction, so it mostly measures fluid. Cystic nodules,
mostly benign, are dark because of their fluid, while TI-RADS judges echogenicity on the solid
part. A seventh feature computed on the non-fluid pixels only (run J) removed the reversed
direction but added no signal; with it the language model was slightly worse (test interval
$[-0.043, -0.001]$ against run I), so the six-feature set is kept. Margin sharpness may measure
contrast rather than blur, and the simple spot detector finds no signal for microcalcifications in
speckled images.

Added to the 16 radiomics features, the six features raise XGBoost from 0.805 to 0.831 on test
(Table~\ref{tab:results-clinical-xgb}), with the same gain on validation and calibration. Adding the
remaining 38 radiomics features instead gives nothing (0.796), so the gain comes from information
that the radiomics features do not contain. The six features alone reach 0.775.

\begin{table}[ht]
  \centering
  \small
  \caption{XGBoost with and without the TI-RADS-motivated features (pooled split, radiomics
           features after the near-duplicate filter).}
  \label{tab:results-clinical-xgb}
  \begin{tabular}{lrrrrr}
    \toprule
    features & number & val & cal & test & \multicolumn{1}{p{2.4cm}}{\raggedleft cancers in lowest 15\% of cal} \\
    \midrule
    radiomics, mRMR 16            & 16 & 0.799 & 0.823 & 0.805 & 4/78 \\
    radiomics, mRMR 16 + TI-RADS  & 22 & 0.828 & 0.848 & 0.831 & 2/78 \\
    TI-RADS only                  &  6 & 0.786 & 0.787 & 0.775 & 6/78 \\
    radiomics, all                & 54 & 0.805 & 0.816 & 0.796 & 2/78 \\
    radiomics, all + TI-RADS      & 60 & 0.824 & 0.836 & 0.826 & 3/78 \\
    \bottomrule
  \end{tabular}
\end{table}

The language model gains more. Because a single run can be off by 0.03--0.05
(Section~\ref{sec:results-seeds}), the comparison was repeated with three matched seeds: each seed
trains the model once with and once without the six features
(Table~\ref{tab:results-clinical-seeds}).

\begin{table}[ht]
  \centering
  \small
  \setlength{\tabcolsep}{4pt}
  \caption{Language model with and without the TI-RADS-motivated features, three matched seeds
           (42, 1, 2). Mean $\pm$ standard deviation over seeds; the last row averages the three
           seeds' probabilities per nodule and gives a 95\% paired bootstrap interval.}
  \label{tab:results-clinical-seeds}
  \begin{tabular}{lrrr}
    \toprule
    & val & cal & test \\
    \midrule
    with TI-RADS (I, M, N)        & $0.785 \pm 0.012$ & $0.794 \pm 0.017$ & $0.785 \pm 0.018$ \\
    without (D, O, P)             & $0.727 \pm 0.022$ & $0.721 \pm 0.029$ & $0.726 \pm 0.024$ \\
    difference per seed, mean (min) & $+0.058$ ($+0.045$) & $+0.074$ ($+0.059$) & $+0.059$ ($+0.050$) \\
    seed-averaged, 95\% interval  & $[+0.018, +0.092]$ & $[+0.033, +0.100]$ & $[+0.019, +0.094]$ \\
    \bottomrule
  \end{tabular}
\end{table}

The gain is positive for every seed and every split, at least $+0.045$, and the test interval of
each matched pair excludes zero (seed 42: $[+0.015, +0.084]$; seed 1: $[+0.013, +0.098]$; seed 2:
$[+0.026, +0.112]$). Seed 42 gave the best run on both sides, so the single-run comparison of
runs I and D (+0.050) did not overstate the gain: on average it is about $+0.06$, slightly more
than for XGBoost. Training with the features was also steadier: validation AUC rose over the
epochs in all three seeds.

\subsection{Peritumoral radiomics}
\label{sec:results-peri}

Since information about the surroundings helped, the PyRadiomics features were also computed in a
ring of tissue around the nodule (width half of the nodule's equivalent radius, at least 5
pixels; all families except shape; 93 features, 49 after the near-duplicate filter;
\texttt{conformal\_triage/peritumoral.py}). The rule was fixed before running: the ring features
would be given to the language model only if adding eight of them (mRMR) to XGBoost improved the
calibration AUC with an interval above zero.

\begin{table}[ht]
  \centering
  \small
  \caption{XGBoost with peritumoral (ring) features, pooled split.}
  \label{tab:results-peri}
  \begin{tabular}{lrrrr}
    \toprule
    features & number & val & cal & test \\
    \midrule
    nodule mRMR 16 + TI-RADS                    & 22 & 0.828 & 0.848 & 0.831 \\
    \quad + ring mRMR 8                         & 30 & 0.832 & 0.851 & 0.834 \\
    \quad + all ring features                   & 71 & 0.831 & 0.853 & 0.832 \\
    joint mRMR 16 (nodule and ring) + TI-RADS   & 22 & 0.836 & 0.840 & 0.821 \\
    ring mRMR 16 only                           & 16 & 0.699 & 0.671 & 0.693 \\
    \bottomrule
  \end{tabular}
\end{table}

Adding the ring features changes the calibration AUC by $+0.002$, with a 95\% interval from
$-0.009$ to $+0.014$, and the test AUC by $+0.003$ (from $-0.010$ to $+0.016$). By the rule they
were not passed to the language model. The ring carries signal on its own (test 0.693), but it is redundant with the nodule and
TI-RADS-motivated features: the ring features mRMR selects are mostly brightness statistics, close
to what the echogenicity ratio already compares. What helped was not more features about the
nodule's neighbourhood but the specific information TI-RADS points to, such as orientation and
outline, that standard radiomics does not measure.

\section{How does the language model compare with XGBoost?}
\label{sec:results-xgb}

XGBoost is better with the same features, by about 0.05 AUC, and the language model adds nothing
to it.

\begin{itemize}
  \item Official split, 8 v1 features: XGBoost 0.696, language model 0.667; with 16 features
        XGBoost reaches 0.709 while the 7B language model stays at 0.666.
  \item Pooled split, 16 v2 features after filtering: XGBoost 0.805, language model 0.754
        (run D).
  \item Pooled split, the same 22 features as run I: XGBoost 0.835, language model
        $0.785 \pm 0.018$ over three seeds.
\end{itemize}

Run I alone (0.804) had suggested that the gap shrinks to 0.03 with the TI-RADS features; over
seeds it stays near 0.05.

\subsection{Fusion of the two models}
\label{sec:results-fusion}

A smaller AUC does not rule out that the language model carries information XGBoost misses. To
test this, XGBoost was trained on exactly the 22 features of run I with the same preprocessing,
and the two models' probabilities were combined (\texttt{experiments/hybrid\_xgb\_llm.py}). The
language model is the average of its three seeds.

\begin{table}[ht]
  \centering
  \caption{XGBoost, the language model and their combinations, on the 22 features of run I.
           Stacking is a logistic regression on the two logits, fitted on validation only.}
  \label{tab:results-fusion}
  \begin{tabular}{lrrr}
    \toprule
    model & val & cal & test \\
    \midrule
    XGBoost                                  & 0.827 & 0.848 & \textbf{0.835} \\
    language model, average of 3 seeds       & 0.792 & 0.803 & 0.793 \\
    mean of the two probabilities            & 0.820 & 0.838 & 0.824 \\
    mean of the two ranks                    & 0.819 & 0.837 & 0.824 \\
    stacking                                 & 0.827 & 0.848 & 0.834 \\
    \bottomrule
  \end{tabular}
\end{table}

\begin{itemize}
  \item The language model is below XGBoost on every split (test difference $-0.042$, interval
        $[-0.067, -0.017]$).
  \item Averaging the two is slightly worse than XGBoost alone (test $-0.011$, interval
        $[-0.023, +0.001]$).
  \item The stacker gives XGBoost a weight of 0.99 and the language model $-0.04$: it learns to
        ignore the language model, and its AUC equals that of XGBoost.
  \item The two models' probabilities correlate 0.87--0.88 on every split. Two seeds of the
        language model correlate 0.91, so the language model differs from XGBoost hardly more
        than from itself.
\end{itemize}

On these features the fine-tuned language model reproduces a weaker version of the tabular
model's ranking and carries no information that XGBoost does not already use.

\section{Can Learn-then-Test certify automatic decisions?}
\label{sec:results-ltt}

No. In none of the runs, with any score, did the planned LTT procedure certify a single automatic
decision at $\alpha_1 = 5\%$ or $\alpha_2 = 20\%$, and neither did it for XGBoost. This section
explains why, starting with the perturbation votes of the original design.

\subsection{Mixed perturbations push the model off its format}

In run 1 (official split, validation nodules), shuffling the feature order alone changed 29\% of
the answers. Paraphrased templates changed 47--48\% of the answers and were \emph{never}
classified malignant, in any split. No nodule received more than 6 of 10 malignant votes, so the
auto-malignant side was impossible by construction. The model also depended on feature position:
shuffling changed only 10\% of the answers when the first feature stayed first, against 24--39\%
otherwise. These perturbations measured the distance from the training format, not the
uncertainty about the nodule.

\subsection{Format-preserving votes are not better than softmax}

With jitter and drop-one perturbations (run 2), the model stays in distribution: noise and
dropping change 5--12\% of the answers and the AUC of the perturbed queries stays within 0.02 of
the base prompt. But the votes separate correct from incorrect predictions no better than the
softmax probability (Table~\ref{tab:results-errsep}). Jitter votes are worse on calibration and
test; drop votes are level with softmax.

\begin{table}[ht]
  \centering
  \caption{AUC for separating correct from incorrect predictions (run 2, official split).}
  \label{tab:results-errsep}
  \begin{tabular}{lrrr}
    \toprule
    confidence score & validation & calibration & test \\
    \midrule
    softmax confidence          & 0.602 & 0.629 & 0.577 \\
    jitter vote agreement (SD 0.2) & 0.560 & 0.579 & 0.533 \\
    drop vote agreement         & 0.620 & 0.652 & 0.571 \\
    \bottomrule
  \end{tabular}
\end{table}

\subsection{The model is consistently wrong}

The design assumed that the model would be unstable on the nodules it gets wrong, so that their
votes would split and they would be referred. The opposite happened. On calibration, 19.4\% of the
nodules with 0 of 10 malignant votes are malignant, and 45.7\% of those with 10 of 10 are benign.
Between 56\% and 69\% of the base prompt's errors are unanimous: all ten queries give the same
wrong answer. Consistency measures how stable an answer is, not whether it is correct, and no vote
threshold can separate these errors.

\subsection{Nothing certified, for any model}

Table~\ref{tab:results-ltt} shows the cleanest auto-benign group available in each setting. All
are far from what LTT needs: 59 nodules without a cancer, 93 with one or 124 with two
(Section~\ref{sec:methods-triage}).

\begin{table}[ht]
  \centering
  \small
  \caption{Cleanest auto-benign groups on calibration (524 nodules, 182 cancers on the pooled
           split).}
  \label{tab:results-ltt}
  \begin{tabular}{p{4.2cm}p{7.4cm}}
    \toprule
    setting & best auto-benign group \\
    \midrule
    official, jitter votes  & $v = 0$: 165 nodules, 19.4\% malignant \\
    official, drop votes    & $v = 0$: 175 nodules, 20.0\% malignant \\
    official, softmax       & not certified at any candidate \\
    pooled, run E           & lowest 15\%: 78 nodules, 1 malignant; $P \le 0.20$: 90 nodules, 3 malignant \\
    pooled, run I           & lowest 15\%: 78 nodules, 8 malignant; $P \le 0.05$: 4 nodules, none malignant \\
    pooled, LLM, mean of 3 seeds & lowest 10\%, 15\%, 20\%: 4/52, 7/78, 10/104 malignant \\
    pooled, XGBoost (22 features) & lowest 10\%, 15\%, 20\%: 2/52, 2/78, 3/104 malignant \\
    \bottomrule
  \end{tabular}
\end{table}

Run E comes closest among the language models: its 78 lowest-scored calibration nodules contain
one cancer (1.3\%), well below the 5\% target, but with one cancer 93 nodules are needed. The best
setting is much dirtier there: 8 cancers among the lowest 78 for run I and 7 for the average of its
three seeds. XGBoost, with the TI-RADS
features, has the cleanest tail of all models (2 cancers among the lowest 78, 3 among the lowest
104), and still nothing is certified: two cancers need 124 nodules, three need 153.

The auto-malignant side shows a second obstacle, which lies in the procedure rather than the
model. For XGBoost, the candidate $u = 0.85$ would pass on its own (54 calibration nodules, 4
benign, $p = 0.010$ against $\alpha_2 = 20\%$). The planned fixed sequence, however, starts at the
strictest threshold, $u = 0.95$, where only 7 nodules fall and 3 of them are benign; it fails
there, stops, and never reaches 0.85. Run I has the same pattern ($u = 0.85$: 50 nodules, 4
benign, $p = 0.018$; $u = 0.95$: 13 nodules, $p = 0.23$). The strictest candidates hold too few
nodules to ever pass, so a sequence that starts with them stops at once.

Ordering the candidates by their $p$-values on the validation split, which is independent of
calibration and therefore keeps the guarantee, avoids starting at a near-empty threshold. This
certifies auto-malignant at $u = 0.85$ for run I (test: 43 nodules, 9.8\%, of which 5 benign,
11.6\%) and at $u = 0.90$ for the language model averaged over three seeds (test: 28 nodules,
6.4\%, none benign). It certifies nothing for runs D and E, nothing for XGBoost (whose validation
order still puts 0.95 first), and nothing on the auto-benign side. Because the ordering rule was
chosen after run I was seen, this result is exploratory: it has to be fixed in advance and
confirmed on another split before it can be claimed.

\subsection{A design fixed in advance, over re-splits}
\label{sec:results-ltt-design}

The two obstacles suggest a design that keeps the guarantee and can pass. It was written down and
committed before it was run; the code is in
\texttt{experiments/\allowbreak ltt\_design\_resplits.py}.

\begin{itemize}
  \item \textbf{Candidates} at quantiles of the validation scores, so that each holds a known
        share of the nodules: auto-benign below the 5\%, 10\%, \dots, 30\% quantile, auto-malignant
        above the 95\%, 90\%, \dots, 70\% quantile.
  \item \textbf{Order} by the candidates' binomial $p$-values on validation, smallest first.
  \item \textbf{Test} on calibration with the fixed sequence, $\alpha_1 = 5\%$,
        $\alpha_2 = 20\%$, $\delta = 0.05$, as planned.
\end{itemize}

Keeping the models fixed, calibration and test were re-split 2,000 times as in
Section~\ref{sec:results-splitcp}, and the new design was compared with the planned one
(Table~\ref{tab:results-ltt-design}).

\begin{table}[ht]
  \centering
  \small
  \setlength{\tabcolsep}{4pt}
  \caption{Share of 2,000 re-splits in which a side is certified, and, when the quantile design
           certifies auto-malignant, the share of test nodules decided and the share of benign
           nodules among them (on test, and on all 962 nodules above the threshold).}
  \label{tab:results-ltt-design}
  \begin{tabular}{lrrrrr}
    \toprule
    & \multicolumn{1}{p{1.5cm}}{\raggedleft auto-benign, either design}
    & \multicolumn{1}{p{1.6cm}}{\raggedleft auto-malig., planned}
    & \multicolumn{1}{p{1.6cm}}{\raggedleft auto-malig., quantile}
    & \multicolumn{1}{p{1.4cm}}{\raggedleft decided, test}
    & \multicolumn{1}{p{1.9cm}}{\raggedleft benign among them (test / all)} \\
    \midrule
    XGBoost (22 features) & 0\% & 0\%  & 38\% & 10.7\% & 14.9\% / 11.2\% \\
    LLM, mean of 3 seeds  & 0\% & 1\%  & \textbf{99\%} & 7.9\% & 7.4\% / 6.4\% \\
    LLM run I             & 0\% & 13\% & 22\% & 11.5\% & 18.4\% / 13.6\% \\
    LLM run D             & 0\% & 0\%  & 8\%  & 6.8\%  & 22.6\% / 13.8\% \\
    \bottomrule
  \end{tabular}
\end{table}

\begin{itemize}
  \item \textbf{Auto-benign at 5\% is never certified}, for any model or design. This side is out
        of reach of the current models; no choice of procedure changes that.
  \item \textbf{Auto-malignant at 20\% becomes attainable.} For the language model averaged over
        its three seeds, the new design certifies in 99\% of re-splits, against 1\% for the
        planned design. About 8\% of the nodules are then called malignant automatically, 6--7\%
        of them benign. On our split the rule is $u = 0.833$: 9.8\% of the test nodules, 14\% of
        them benign.
  \item \textbf{The order only helps where validation and calibration agree.} For XGBoost the
        validation order puts its top 5\% first, which on calibration contains 3 benign nodules
        among 7, so the sequence often stops there (38\%). The seed-averaged language model, whose
        scores are smoother, benefits most.
  \item \textbf{The bound holds.} The share of benign nodules among all 962 nodules above the
        chosen threshold never exceeds 20\% in any re-split. The test-only share can (run D):
        in a fixed pool, a re-split with a clean calibration part leaves a dirtier test part, so
        conditioning on certification biases the test share upwards.
\end{itemize}

Two caveats limit this result. The design was written after the exploratory result on run I, and
the models were selected on these data; the re-splits vary only which nodules calibrate and which
test. Confirmation needs new data: retraining on new splits, or a second dataset.

\subsection{A guarantee on missed cancers}

Bounding the share of \emph{all} cancers sent to auto-benign (Section~\ref{sec:methods-miss}) is
easier, because the evidence comes from all 182 calibration cancers. At 5\% it is certified in
every pooled run (Table~\ref{tab:results-miss}).

\begin{table}[ht]
  \centering
  \small
  \caption{Miss-rate guarantee at 5\%, pooled split. LLM avg.: the language model averaged over
           runs I, M and N; XGB: XGBoost on the 22 features of run I.}
  \label{tab:results-miss}
  \begin{tabular}{lrrrr}
    \toprule
    run & threshold & auto-benign, cal & auto-benign, test & cancers missed, test \\
    \midrule
    A & 0.16 & 14.3\% & 14.4\% & 5.3\% \\
    B & 0.22 &  4.0\% &  5.7\% & 2.6\% \\
    C & 0.09 & 10.3\% & 11.4\% & 2.6\% \\
    D & 0.16 & 16.8\% & 18.0\% & 8.6\% \\
    E & 0.21 & 19.1\% & 19.9\% & 5.9\% \\
    F & 0.26 &  9.2\% & 11.0\% & 3.3\% \\
    G & 0.17 &  6.9\% &  7.8\% & 0.7\% \\
    H & 0.20 &  4.2\% &  3.7\% & 0.7\% \\
    I & 0.09 &  6.1\% &  5.5\% & 0.7\% \\
    M & 0.09 &  5.2\% &  5.3\% & 1.3\% \\
    N & 0.12 &  9.5\% &  8.4\% & 2.6\% \\
    \midrule
    LLM avg. & 0.14 & 11.3\% & 12.1\% & 4.6\% \\
    XGB & 0.14 & 19.8\% & 23.7\% & 4.6\% \\
    \bottomrule
  \end{tabular}
\end{table}

The realised test miss rate ranges from 0.7\% to 8.6\%. A guarantee with $\delta = 0.05$ is
allowed to fail on some test sets, and the test set has only 152 cancers, so one cancer is 0.7
percentage points; one run above 5\% is not a contradiction. On the official split the same rule
behaved worse under the shift: at 10\%, 15.3\% of the official-test cancers were missed. XGBoost
sends twice as many nodules to auto-benign as the averaged language model at the same miss rate,
and the cancers among them are 7 of 104 (6.7\%) against 7 of 53 (13.2\%).

\subsection{Split conformal prediction on the pooled split}
\label{sec:results-splitcp}

Classical split conformal prediction (Section~\ref{sec:methods-splitcp}) gives what it promises on
the pooled split: with marginal $\alpha = 0.10$ the test coverage is 0.88--0.93, and with
$\alpha = 0.05$ it is 0.93--0.97. Marginal coverage, however, is carried by the benign majority:
at $\alpha = 0.10$ only 66--92\% of the test cancers have ``malignant'' in their set. The Mondrian
form protects the cancers directly (Table~\ref{tab:results-splitcp}).

\begin{table}[ht]
  \centering
  \small
  \caption{Mondrian split conformal prediction on the pooled test set ($\alpha = 0.10$ benign,
           $0.05$ malignant; thresholds from calibration). Auto-benign and auto-malignant are
           the single-label sets; the rest is referred. XGB: XGBoost on the 22 features of run I.}
  \label{tab:results-splitcp}
  \begin{tabular}{lrrrrrr}
    \toprule
    run & \multicolumn{1}{p{1.6cm}}{\raggedleft cancers covered}
        & \multicolumn{1}{p{1.3cm}}{\raggedleft auto-benign}
        & \multicolumn{1}{p{1.7cm}}{\raggedleft cancers in auto-benign}
        & \multicolumn{1}{p{1.5cm}}{\raggedleft auto-malignant}
        & \multicolumn{1}{p{1.7cm}}{\raggedleft benign in auto-malignant}
        & referred \\
    \midrule
    A & 0.901 & 22.8\% & 15.0\% & 19.6\% & 30.2\% & 57.5\% \\
    B & 0.914 & 17.1\% & 17.3\% & 17.8\% & 34.6\% & 65.1\% \\
    C & 0.941 & 18.7\% & 11.0\% & 16.4\% & 41.7\% & 64.8\% \\
    D & 0.908 & 21.7\% & 14.7\% & 17.4\% & 30.3\% & 61.0\% \\
    E & 0.934 & 22.8\% & 10.0\% & 16.7\% & 27.4\% & 60.5\% \\
    F & 0.928 & 18.5\% & 13.6\% & 15.8\% & 39.1\% & 65.8\% \\
    G & 0.954 & 18.0\% &  8.9\% & 14.6\% & 35.9\% & 67.4\% \\
    H & 0.987 &  8.7\% &  5.3\% & 12.3\% & 27.8\% & 79.0\% \\
    I & 0.941 & 18.3\% & 11.3\% & 22.8\% & 22.0\% & 58.9\% \\
    J & 0.914 & 24.7\% & 12.0\% & 22.4\% & 27.6\% & 53.0\% \\
    K & 0.954 & 16.2\% &  9.9\% & 21.5\% & 21.3\% & 62.3\% \\
    L & 0.908 & 14.6\% & 21.9\% & 18.5\% & 33.3\% & 66.9\% \\
    M & 0.914 & 19.6\% & 15.1\% & 19.2\% & 26.2\% & 61.2\% \\
    N & 0.928 & 21.9\% & 11.5\% & 20.3\% & 23.6\% & 57.8\% \\
    O & 0.967 & 12.8\% &  8.9\% & 13.0\% & 40.4\% & 74.2\% \\
    P & 0.941 & 12.8\% & 16.1\% & 15.5\% & 29.4\% & 71.7\% \\
    \midrule
    XGB & 0.941 & 28.1\% & 7.3\% & 24.2\% & 24.5\% & 47.7\% \\
    \bottomrule
  \end{tabular}
\end{table}

About a fifth of the nodules receive the single label ``benign'' (more than a quarter with
XGBoost), and roughly 60\% are referred. Two observations follow.

First, the cancers among the auto-benign calls make up 5--22\% of that group: the guarantee limits
how many cancers are missed overall, not how safe a single benign call is, so it does not replace
$R_1$.

Second, the share of cancers covered is below the 0.95 target in 12 of the 16 runs, which looks
like a systematic failure. All runs, however, share one calibration set and one test set, so this
is one draw of the data seen many times. To separate chance from a real shortfall, the model was
held fixed and the 962 calibration and test nodules were randomly re-split 4,000 times, with the
same sizes and class proportions, recalibrating each time (\texttt{experiments/conformal\_checks.py};
Table~\ref{tab:results-resplit}).

\begin{table}[ht]
  \centering
  \small
  \caption{Mondrian split conformal prediction over 4,000 random re-splits of calibration and
           test (model fixed). Mean and standard deviation of the cancer coverage; mean share of
           auto-benign nodules and of cancers among them, with the 5th--95th percentile range.}
  \label{tab:results-resplit}
  \setlength{\tabcolsep}{4pt}
  \begin{tabular}{lrrrr}
    \toprule
    model & \multicolumn{1}{p{1.8cm}}{\raggedleft coverage, our split}
          & \multicolumn{1}{p{2.1cm}}{\raggedleft coverage, re-splits}
          & auto-benign
          & \multicolumn{1}{p{2.5cm}}{\raggedleft cancers in auto-benign} \\
    \midrule
    XGBoost (22 features)        & 0.941 & $0.951 \pm 0.024$ & 25.7\% & 6.4\% (2.3--10.5\%) \\
    LLM, mean of 3 seeds         & 0.934 & $0.950 \pm 0.024$ & 16.7\% & 9.8\% (4.7--14.6\%) \\
    run I                        & 0.941 & $0.951 \pm 0.024$ & 16.9\% & 9.6\% (4.1--14.0\%) \\
    run D                        & 0.908 & $0.951 \pm 0.023$ & 17.6\% & 9.3\% (3.6--14.9\%) \\
    \bottomrule
  \end{tabular}
\end{table}

The average coverage is 0.950--0.951 for every model, exactly as the theory promises, and falls
below 0.95 in about 45\% of the re-splits. Our split was unlucky by about one standard deviation;
the guarantee holds. The same analysis gives the expected share of cancers among the auto-benign
calls: about 10\% for the language model and 6\% for XGBoost, against 35\% in the data and the 5\%
asked of $R_1$.
